\subsection{Тепловые двигатели}
%%%%%%%%%%%%%%%%%%%%%%%%%%%%%%%%%%%%%%%%%%%%%%%%%%%%
%%%%%%%%%%%%%%%%%%%%%%%%%%%%%%%%%%%%%%%%%%%%%%%%%%%%
 \z{Автомобиль прошел $80$ км. Двигатель автомобиля развивает мощность $P = 40$ кВт и израсходовал $V = 14$ л бензина. С какой средней скоростью двигался автомобиль, если КПД его двигателя $\eta = 30\%$?}
%
{$v = \frac{SP}{qV\rho\eta} = 23.4\,\text{м/с} = 84.2\,\text{км/ч}$}

%%%%%%%%%%%%%%%%%%%%%%%%%%%%%%%%%%%%%%%%%%%%%%%%%%%%
 \z{В подъемнике используется двигатель внутреннего сгорания с КПД $\eta = 30\%$. С помощью подъемника песок массой $m = 150$ т переместили из карьера глубиной $h = 40$ м на поверхность земли. Сколько литров бензина израсходовали?}
%
{$V = \frac{MgH}{\rho\eta q} = 6\,\text{л}$}

%%%%%%%%%%%%%%%%%%%%%%%%%%%%%%%%%%%%%%%%%%%%%%%%%%%%
 \z{С какой постоянной скоростью ехал поезд, если за $t = 30$ мин в печи его двигателя сгорело $m = 18$ кг древесного угля? КПД двигателя равен $\eta = 30\%$, а его сила тяги $F = 5$ кН. Удельная теплота сгорания древесного угля $q = 34$ МДж/кг.}
%
{$v = \frac{mq\eta}{Ft} = 20.4\,\text{м/с} = 73.4\,\text{км/ч}$}

%%%%%%%%%%%%%%%%%%%%%%%%%%%%%%%%%%%%%%%%%%%%%%%%%%%%
 \z{Как известно, лошадь может развивать мощность, равную одной лошадиной силе. Одна лошадиная сила примерно равна $178$ кал/с. Калорийность травы $10$ Ккал/$100$ г. Сколько травы за час должна съедать работающая лошадь, чтобы поддерживать себя в форме? Считать, что КПД усвоения травы $\beta = 80\%$. Ответ выразить в кг, округлить до десятых.}
%
{$\mu = \frac{Pt}{\beta q} = 8\,\text{кг/час}$}

%%%%%%%%%%%%%%%%%%%%%%%%%%%%%%%%%%%%%%%%%%%%%%%%%%%%
 \z{Идеальная тепловая машина получила от нагревателя $Q_{\text{н}}$. Сколько тепла она отдаст холодильнику? Температура нагревателя $T_{\text{н}}$, холодильника – $T_{\text{х}}$.}
%
{$Q_{\text{х}} = Q_{\text{н}} \cdot \frac{T_{\text{х}}}{T_{\text{н}}}$}

%%%%%%%%%%%%%%%%%%%%%%%%%%%%%%%%%%%%%%%%%%%%%%%%%%%%
 \z{Идеальная тепловая машина совершила работу $A$. Сколько тепла она отдаст холодильнику? Температура нагревателя $T_{\text{н}}$, холодильника – $T_{\text{х}}$.}
%
{$Q_{\text{х}} = A \cdot \frac{T_{\text{н}}}{T_{\text{н}}-T_{\text{х}}}$}

%%%%%%%%%%%%%%%%%%%%%%%%%%%%%%%%%%%%%%%%%%%%%%%%%%%%
 \z{Сколько воды, находящейся при температуре $100\degree\text{C}$, можно испарить, имея в распоряжении $1$ кг воды при температуре $0\degree\text{C}$?}
%
{$m_{\text{п}} = m_{\text{л}} \cdot \frac{\lambda}{L} \cdot \frac{T_{100}}{T_0} = 200\,\text{г}$}